\documentclass[11pt]{article}
\usepackage[T1]{fontenc}
\usepackage{lmodern}
\usepackage{amsmath,amssymb,booktabs}
\usepackage[margin=1in]{geometry}
\usepackage[hidelinks]{hyperref}
\title{When Do Smoothed Financial Risk Indicators Fail to Forecast?\\
Evidence Across Countries and Financial Disturbances}
\author{Working draft --- authorship to be confirmed}
\date{16 October 2026}
\begin{document}
\maketitle
\begin{abstract}
We set out a comparative empirical design to study the historical reliability
of smoothed-trend forecasts for country-specific financial-risk indices.
Risk measurement and trend extrapolation are treated as separate statistical
operations. We propose transparent signed aggregates, principal-component
indices and dynamic-factor alternatives, followed by a Gaussian
marginal-likelihood-calibrated trend smoother. Both full-sample historical
description and strictly chronological forecast evaluation will be considered,
with the latter providing the evidence for predictive claims. The study will
compare predetermined crisis and non-crisis periods and report interval
calibration and conventional forecast benchmarks. This document is a
prospective design; new empirical results are not yet available.
\end{abstract}
\section{Research question}
In financial disturbances, the trajectories implied by a recently estimated
smooth trend may diverge markedly from later observations. It remains unclear
whether such divergence indicates informative departures from prevailing
conditions or merely ordinary shifts in forecast uncertainty.
Our question is: across countries, index-construction methods, and forecast
horizons, when do smoothed financial-stress-index forecasts cease to be
reliable, and do errors reveal information beyond conventional models?

\subsection{Scope and contribution}
The intended contribution is a transparent cross-country empirical assessment
rather than a new penalized smoothing formula, numerical optimizer or
dynamic-factor estimator. We plan to study Mexico, the United States,
the United Kingdom, Japan and a defined euro-area aggregate, subject to
audited series availability. Earlier U.S.-only retrospective results motivate,
but do not establish, the proposition.

\section{Measurement and forecasting methods}
\subsection{Risk index construction}
Let \(X_{c,t}\) denote a country-specific vector of
transformed financial indicators, each oriented and scaled under
a declared convention. Candidate indices include an economically
signed composite, the first principal component, and a latent state
\(f_{c,t}\) from a dynamic factor representation
\begin{equation}
X_{c,t}=\Lambda_c f_{c,t}+u_{c,t}.
\end{equation}
The choice of raw series, transformations, anchor indicators and
historical availability dates belongs to the measurement stage.

\subsection{Likelihood-calibrated trend}
Conditionally on an estimated index \(y_t\), the primary proposed
Gaussian specification is
\begin{align}
y_t &= \tau_t+\varepsilon_t,
&\varepsilon_t&\sim N(0,r),\\
\Delta^2\tau_t &= \eta_t,
&\eta_t&\sim N(0,q),\qquad r,q>0.
\end{align}
With compatible diffuse treatment of the linear-trend nullspace,
the corresponding penalized least-squares criterion is
\begin{equation}
\widehat{\tau}_{\lambda} =
\arg\min_{\tau}
\left\{ \|y-\tau\|^2+
\lambda\|D_2\tau\|^2\right\},
\qquad \lambda=r/q.
\end{equation}
The two positive variances will be estimated by maximizing
the \emph{marginal observed-data likelihood}, e.g. through the Kalman
filter. No such implementation has yet been validated in this project.
The trend-only transition specification is deliberately narrower than
a general local-linear-trend model with an additional level innovation.

\subsection{Forecast error}
For an explicitly defined estimated-factor target at horizon \(h\),
\begin{equation}
e_{c,t,h}^{(m)} =
f_{c,t+h}^{(m),\mathrm{target}} -
\widehat f_{c,t+h\mid t}^{(m)},
\quad
z_{c,t,h}^{(m)}=
\frac{e_{c,t,h}^{(m)}}{\widehat\sigma_{c,t,h}^{(m)}}.
\end{equation}
Large forecast errors may signal misspecified dynamics, heavy tails,
variance changes, unreliable factor measurement, or structural change.
They must not be treated automatically as crisis evidence.

\section{Prospective evaluation}
\subsection{Retrospective description versus pseudo-real-time analysis}
A full-sample historical factor provides a useful retrospective view
of financial disturbances, but its parameters and normalization may
use information from dates after the purported forecast origin.
Consequently, the primary forecasting experiment will instead use
historically available data and refit transformations, the index, the
trend model, and forecast calibration at each origin.
Retrospective and pseudo-real-time estimates will be labeled separately.

\subsection{Countries, events and benchmark forecasting methods}
Proposed countries comprise Mexico, the United States, the United
Kingdom, Japan and a euro-area aggregate. Predeclared case studies may
include the 2008--2009 financial crisis, European sovereign debt stress,
the COVID-19 shock, and later monetary and banking disturbances.
The 1994--1995 Mexican crisis is a separate possible extension because
the proposed contemporary indicators have limited older coverage.
All event and quiet-control windows will be fixed before confirmatory tests.

\subsection{Statistical evaluation}
For each index, assess matched-origin forecast RMSE, MAE, predictive
interval coverage and standardized exceedance rates against a random walk,
an autoregression, fixed trend smoothness, and (where meaningful) a
factor-transition forecast. Cross-index direct performance comparisons
require a predeclared \emph{common observable target}: indices are not
assumed to be interchangeable random variables. Dependence across
overlapping forecast horizons and shared country shocks must enter the
uncertainty analysis. The complete prospective protocol is EXP-0001.

\section{Results}
No cross-country empirical results have yet been computed or validated.
This section will be completed only after reproducible experiments are run.
\section{Discussion and limitations}
An estimated factor is not directly observed financial risk. Factor sign and
scale are model conventions; alternative index constructions yield distinct
targets that should not be ranked by unadjusted cross-model factor RMSE.
Predictive intervals require an uncertainty model that addresses heavy tails,
heteroskedasticity and estimation uncertainty.
\section{Conclusions}
Conclusions will be written after the experimental protocol and analysis have
been completed. The existence of forecast errors is not in itself proof of
structural change or early-warning value.
\begin{thebibliography}{9}
\bibitem{guerrero2007}
V.~M. Guerrero (2007).
Time series smoothing by penalized least squares.
\emph{Statistics \& Probability Letters}, 77(12), 1225--1234.
\bibitem{islas2019}
A. Islas, V.~M. Guerrero, and E. Silva (2019).
Forecasting remittances to Mexico with a multi-state Markov-switching model
applied to the trend with controlled smoothness.
\emph{Romanian Journal of Economic Forecasting}, 22(1).
\bibitem{biessy2026}
G. Biessy (2026).
Whittaker--Henderson smoothing revisited: A modern statistical framework
for practical use.
\emph{ASTIN Bulletin}. doi:10.1017/asb.2025.10061.
\end{thebibliography}
\end{document}
